\documentclass[11pt]{article}
\usepackage[a4paper,margin=27mm]{geometry}
\usepackage{amsmath,amssymb,amsthm,hyperref}
\setlength{\emergencystretch}{2em}
\newtheorem{theorem}{Theorem}
\newtheorem{lemma}{Lemma}
\title{Realizing scalar comparison profiles by finite uniform dice}
\author{Research proof; independent mathematical review completed}
\date{30 September 2026}
\begin{document}\maketitle
No bound on the number of faces is asserted in this note. The point is
to transfer robust local comparison profiles from a scalar quadrature
measure to a completely finite, exactly permutation-fair family.

\begin{theorem}\label{main}
Let $m\ge2$ and let $\sigma$ be a probability measure on $[0,1]$
whose moments through degree $m$ equal those of Lebesgue measure.
Fix $t\in(0,1)$ with $\sigma(\{t\})=0$, and let $\rho>0$.
There is a permutation-fair family of $m+1$ finite uniform dice,
with pairwise distinct face labels, and a face of one die $D$ at which
the probabilities $p_i$ that the other dice fall below that face satisfy
\[
 |p_1-F_\sigma(t)|<\rho,\qquad
 |p_i-t|<\rho\quad(2\le i\le m).
\]
In particular the weights of all rows of the comparison tensor of $D$
are equal. The remaining $m$ dice can be partitioned into one special
die and $m-1$ dice with the same number of faces as $D$.
\end{theorem}

\begin{lemma}\label{approx}
Under the hypotheses of Theorem~\ref{main}, for every $\eta>0$ there
is a degree-$m$ equal-weight quadrature for Lebesgue measure, with
$K\ge m^2$ distinct nodes in $(0,1)$, none equal to $t$, and empirical
measure $\nu$ satisfying $|F_\nu(t)-F_\sigma(t)|<\eta$.
\end{lemma}
\begin{proof}
We first approximate $\sigma$ by a bounded, strictly positive density
with exactly the same first $m$ moments. Smooth each point of $\sigma$
on a small interval contained in $(0,1)$, including points near the
endpoints, to obtain a bounded density $g_a$ converging weakly to
$\sigma$ as $a\downarrow0$. For example, push the measure into
$[a,1-a]$ and convolve with a uniform kernel of radius $a/2$.
Fix $\delta>0$ and form $(1-\delta)g_a+\delta$.
Its moment error against any fixed polynomial tends to zero as
$a\downarrow0$.

Let $\phi_0=1,\phi_1,\ldots,\phi_m$ be the orthonormal shifted
Legendre polynomials. Subtract the polynomial
\[
 h_a=\sum_{j=1}^m\phi_j
       \int_0^1\phi_j(x)((1-\delta)g_a(x)+\delta)\,dx.
\]
For fixed $m,\delta$, its supremum norm tends to zero. Choose $a$
so small that $\|h_a\|_\infty<\delta/2$. Then
\[
 w=(1-\delta)g_a+\delta-h_a
\]
is a bounded density, bounded below by $\delta/2$, and has the
required exact moments. By first making $\delta$ small and then $a$
small, $w(x)\,dx$ can be made arbitrarily close to $\sigma$ weakly.

A density bounded above and below by positive constants is a doubling
weight. The equal-weight quadrature theorem of Gilboa--Peled
(\href{https://arxiv.org/abs/1507.01505}{Theorem 1.1}) supplies
quadratures $\nu_D$ for $w$ of arbitrarily high degree $D$, with
arbitrarily large admissible cardinality. As $D\to\infty$, these
empirical measures converge weakly to $w(x)\,dx$: approximate a
continuous test function uniformly by a polynomial, and use
exactness. For $D\ge m$, each $\nu_D$ is also exact through degree
$m$ for Lebesgue measure.

Choose $D\ge2m+2$ and cardinality $K\ge m^2$. Such a rule has at
least $m+2$ distinct support points; otherwise the square of its
support-vanishing polynomial contradicts exactness and positivity
of $w$. At least $m$ support points are interior. Select $m$ such
nodes as pivot coordinates. The first $m$ moment equations have
an invertible Vandermonde Jacobian in the pivots. The implicit
function theorem permits arbitrarily small perturbations of all
other coordinates, preserving the first $m$ moments by compensating
with the pivots. Endpoints can be moved inward, and all collisions
can be avoided, as in the regularization lemma in
\path{size_bounds/polynomial_exceptional_die.tex}.

We may also avoid the fixed value $t$. For a pivot initially equal
to $t$, vary a free coordinate at a support point different from all
pivots; there is such a point because the original rule has at least
$m+2$ distinct support points. The corresponding pivot derivative
is minus a nonzero Lagrange cardinal value. Thus the constraint
that this pivot equal $t$ is a proper analytic zero set. The same
is immediate for free coordinates. A finite union of these zero
sets and the collision sets cannot cover the permitted open set
of inward perturbations.

These regularizations can be arbitrarily small, so they preserve
weak approximation. Choose the preceding parameters successively
so that the resulting empirical measure is close enough to $\sigma$
at the continuity point $t$. More explicitly, first choose two
continuity points $t_-<t<t_+$ of $\sigma$ whose distribution values
differ by less than $\eta/4$. Weak approximation makes the
distribution values of the quadrature at both endpoints close to
those of $\sigma$; a perturbation of size less than half their
distance to $t$ brackets the new value at $t$ between these values.
This proves the lemma.
\end{proof}

\begin{proof}[Proof of Theorem~\ref{main}]
Apply Lemma~\ref{approx} with error smaller than $\rho/4$, and write
its nodes as $u_1<\cdots<u_K$. The one-exceptional-die construction
in \path{size_bounds/polynomial_exceptional_die.tex} applies to any
such regular rule with $K\ge m^2$, not just to its initial
$O(m^2)$-node choice. Its proof has the following useful additional
freedom: for every $\gamma>0$, the nodes can be replaced by rational
$v_q$ arbitrarily close to $u_q$, and the $m$ old continuous dice
can be given positive rational step densities $f_j$ such that
\[
 \|f_j-1\|_\infty<\gamma\quad(1\le j\le m),
\]
while their joint order law with the uniform special die on the
$v_q$ is exactly fair. Indeed, the fixed separated bump functions
are bounded, and their rational Vandermonde coefficients tend to
zero as $v\to u$. Choose the rational approximation close enough
that no node crosses $t$.

The final finite realization of that construction partitions the
interval into rational constant-density cells and replaces every
cell by an equal-sided fair block with rational mass multiplicities.
It permits any further rational refinement of the partition.
Take a refinement of mesh at most $\tau$, and select a cell $J$
within distance $\tau$ of $t$, contained in the same gap between
special nodes as $t$. Such a positive-length cell exists even if
$t$ is already a density breakpoint. All old densities are positive,
so every old die, including a selected die $D$, has faces in $J$.
Place the discrete block's faces at distinct real coordinates inside
their cells before replacing the total order by integer labels.

Let $x$ be a face of $D$ in $J$. Each other old die has the correct
mass in every cell. Its finite CDF differs from its continuous CDF
at any point by at most the mass of that point's cell, which is at
most $(1+\gamma)\tau$. Its continuous CDF differs from the identity
by at most $\gamma$. Therefore
\[
 |p_i-t|\le\gamma+(2+\gamma)\tau\qquad(2\le i\le m).
\]
The special die has no node between $x$ and $t$, so
$p_1=F_\nu(t)$. Choose $\gamma,\tau$ small enough. This proves all
assertions, including exact fairness and uniform face weights.
\end{proof}

\paragraph{Limit of the statement.}
The approximation parameters, the doubling constant, the quadrature
cardinality, and the final denominator clearing may depend arbitrarily
badly on $m$ and $\rho$. This is a realization theorem for robust
inequalities, not an efficient construction.
\end{document}
